% --- scholar LaTeX fragment --------------------------------------------
% Generated by: scholar sessions export -f table
% Session: classes-missuppfattningar-v2
% \input this file from the body of a document whose
% preamble loads:
%   \usepackage{longtable}
%   \usepackage{booktabs}
%   \usepackage{array}
% ----------------------------------------------------------------------

\begingroup\footnotesize
\begin{longtable}{@{}%
  >{\raggedright\arraybackslash}p{0.44\textwidth}c%
  >{\raggedright\arraybackslash}p{0.13\textwidth}%
  >{\raggedright\arraybackslash}p{0.26\textwidth}@{}}
\caption{Träffar som rör påståendet, missuppfattningar om klasser och objekt (6 citerade, 9 som stöder påståendet och 4 som kvalificerar eller motsäger det, av 387 arbeten; de 152 angränsande och 216 felträffarna listas inte).  Ett arbete som flera databaser lämnat står på en rad, med databaserna uppräknade. Skälet till varje inkludering står i sista kolumnen; för maskinklassade rader anges modellens konfidens. Databaser: OpenAlex (OA; inte open access), DBLP, Web of Science (WoS), Scopus.}\label{tab:sok-missuppfattningar}\\
\toprule
Titel & År & Leverantör & Skäl \\
\midrule
\endfirsthead
\multicolumn{4}{@{}l}{\emph{\tablename~\thetable{} (forts.)}}\\
\toprule
Titel & År & Leverantör & Skäl \\
\midrule
\endhead
\midrule \multicolumn{4}{r@{}}{\emph{forts.\ på nästa sida}}\\
\endfoot
\bottomrule
\endlastfoot
A comparison of the comprehension of object-oriented and procedural programs by novice programmers & 1999 & WoS, OA & \emph{citerad}: programförståelse (oo-vs-procedural; program-comprehension) \\
Avoiding object misconceptions. & 1997 & DBLP & \emph{citerad}: attributen som identitet (object-misconceptions; teaching) \\
Checklists for grading object-oriented CS1 programs: concepts and misconceptions. & 2007 & OA, DBLP, WoS & \emph{citerad}: objektbegreppet (class-object-misconceptions; cs1) \\
Novice Java programmers' conceptions of \textquotedbl{}object\textquotedbl{} and \textquotedbl{}class\textquotedbl{}, and variation theory & 2005 & OA, Scopus & \emph{citerad}: klass som mall eller samling (object-class-conceptions; phenomenography) \\
Object-Oriented Design and Programming: An Investigation of Novices' Conceptions on Objects and Classes & 2015 & WoS & \emph{citerad}: klass som mall eller samling (object-class-conceptions; undergraduates) \\
The Impact of Different Teaching Approaches and Languages on Student Learning of Introductory Programming Concepts & 2016 & OA & \emph{citerad}: undervisningsordning (teaching-approaches; language-comparison) \\
A Diagnostic Tool for Assessing Students' Perceptions and Misconceptions Regards the Current Object \textquotedbl{}this\textquotedbl{}. & 2018 & DBLP & stöder påståendet (0.93) (this-reference; diagnostic-tool) \\
Exploring the Difficulties of Learning Object-Oriented Techniques & 1997 & OA & stöder påståendet (0.95) (basic-object-concepts; learning-difficulties) \\
Object-Oriented Programming: Diagnosis Understanding by Identifying and Describing Novice Perceptions & 2020 & WoS, Scopus & stöder påståendet (0.90) (misconceptions; concept-inventory) \\
On the Students' Misconceptions in Object-Oriented Language Constructs. & 2019 & DBLP & stöder påståendet (0.87) (language-constructs; misconceptions) \\
Statistical frequency-analysis of misconceptions in object-oriented-programming - Regularized PCR models for frequency analysis across OOP concepts and related factors & 2018 & Scopus & stöder påståendet (0.81) (misconception-frequency; oop-concepts) \\
Students' Conceptions of Object-Oriented Programming in the Context of Game Designing in Computing Education: Design of a Ph.D. Research Project to Explore Students' Conceptions in a Long-Term Study & 2021 & Scopus & stöder påståendet (0.80) (students-conceptions; game-design-context) \\
Students' Conceptions of Programming in the Context of Game Design & 2022 & Scopus & stöder påståendet (0.79) (students-conceptions; game-design-context) \\
Thinking in imperative or objects? A study on how novice programmer thinks when it comes to designing an application & 2019 & Scopus & stöder påståendet (0.72) (imperative-vs-objects-thinking; design-thinking) \\
Understanding the \textquotedbl{}this\textquotedbl{} reference in object oriented programming: Misconceptions, conceptions, and teaching recommendations. & 2021 & DBLP & stöder påståendet (0.96) (this-reference; misconceptions) \\
An empirical study of novice program comprehension in the imperative and object-oriented styles & 1997 & OA & kvalificerar eller motsäger påståendet (0.95) (oo-vs-imperative; novice-comprehension) \\
An examination of procedural and object-oriented systems analysis methods: Does prior experience help or hinder performance? & 1999 & OA, WoS & kvalificerar eller motsäger påståendet (0.84) (oo-vs-procedural; novice-analysts) \\
Comparison of Two Approaches when Teaching Object-Orientated Programming to Novices & 2011 & WoS & kvalificerar eller motsäger påståendet (0.98) (objects-first; teaching-order-comparison) \\
Is there any difference in novice comprehension of a small program written in the event-driven and object-oriented styles? & 2002 & Scopus & kvalificerar eller motsäger påståendet (0.72) (oo-vs-event-driven; novice-comprehension) \\
\end{longtable}
\endgroup
